\chapter{Signal optimisation}
\label{chap:optimisation}

\section{Nested problem structure}
Optimisation and traffic assignment solve different problems.  Assignment redistributes
fixed demand for one fixed signal plan.  Optimisation compares alternative signal
plans, each of which requires its own traffic assignment.

\begin{verbatim}
signal candidate
  -> route/lane assignment
  -> dynamic capacity update
  -> lane performance
  -> feasibility and objective
  -> next candidate
\end{verbatim}

\section{Sequential search}
For sequential optimisation, the supplied controller must be the centre of the search.
Candidate zero is the current controller plan:
\begin{verbatim}
cycle  = current cycle
splits = current phase durations
offset = current offset
\end{verbatim}
The current plan is therefore always evaluated before alternatives.  Nearby candidates
are created by transferring the configured split increment between phases while
preserving cycle constraints and minimum greens.  This is preferable to taking the
first combinations from a large arbitrary enumeration because realistic multi-phase
controllers can have many thousands of combinatorial split possibilities.

\section{Minimum green and cycle reconciliation}
A candidate is timing-valid only if every phase respects the configured minimum green
and the controller cycle reconciles with phase greens and intergreens.  For a simple
non-overlap cycle,
\begin{equation}
\sum_p G_p + \sum_p I_p = C.
\end{equation}
The optimiser should not generate candidates that violate its own timing constraints.

\section{Candidate evaluation}
For each candidate, PySig performs route assignment and then evaluates the configured
objective, normally network/intersection delay.  Candidate diagnostics should identify
cycle, phase splits, offset, delay, convergence status and feasibility.

PySig should distinguish hard invalidity (for example, an impossible signal plan,
numerical failure or conservation failure) from poor operational performance.  High
DoS, residual queue and short-lane storage usage are performance measures that should
remain visible to the optimisation objective rather than being silently hidden.

\section{Search limits}
\code{max\_candidates\_per\_controller} is a computational limit.  When the candidate
space is larger than this limit, candidate ordering becomes material.  Starting from
the current plan and exploring nearby candidates gives this limit an operationally
meaningful interpretation.

\section{Convergence}
The optimisation outer loop should consider stability of both signal plans and assigned
flows.  With assignment-responsive RR67 and nested short-lane capacity, changes in
assigned flow can also alter lane capacity.  Capacity-change metrics are therefore
useful additional diagnostics even where the formal stopping criterion remains based on
flow change and signal-plan stability.
